\documentclass[11pt]{article}
\usepackage{amsmath, amssymb}
\usepackage[margin=1in]{geometry}

\begin{document}

\begin{center}
    \textbf{\Large Chapter 2 (Part 2) Notes}\\
    \vspace{0.2cm}
    Tashfeen Hatim Said Omran\\
    UCF ID: 5862894\\
    Class: STA4321
\end{center}

\vspace{0.5cm}

\textbf{\large Page 5 of Ch 2-part 2 notes: Find $P(A \cup B \cup C)$ by first grouping B, C}

I apply the additive law for two events, treating $(B \cup C)$ as a single event:
\[
P(A \cup (B \cup C)) = P(A) + P(B \cup C) - P(A \cap (B \cup C))
\]
I expand $P(B \cup C)$ and apply the distributive law to the intersection:
\[
= P(A) + P(B) + P(C) - P(B \cap C) - P((A \cap B) \cup (A \cap C))
\]
I apply the additive law to $P((A \cap B) \cup (A \cap C))$:
\[
= P(A) + P(B) + P(C) - P(B \cap C) - [P(A \cap B) + P(A \cap C) - P((A \cap B) \cap (A \cap C))]
\]
I simplify the intersection of the intersections to $P(A \cap B \cap C)$ and distribute the negative sign:
\[
P(A \cup B \cup C) = P(A) + P(B) + P(C) - P(A \cap B) - P(A \cap C) - P(B \cap C) + P(A \cap B \cap C)
\]

\vspace{0.5cm}

\textbf{\large Page 5 of Ch 2-part 2 notes: Prove the K-event formula at the bottom of the page using the grouping method}

I use the multiplicative law of probability, $P(X \cap Y) = P(X)P(Y \mid X)$, and group the first $k-1$ events:
\[
P(A_1 \cap A_2 \cap \dots \cap A_k) = P((A_1 \cap A_2 \cap \dots \cap A_{k-1}) \cap A_k)
\]
\[
= P(A_1 \cap A_2 \cap \dots \cap A_{k-1}) P(A_k \mid A_1 \cap A_2 \cap \dots \cap A_{k-1})
\]
I apply the same grouping method to the first term, separating $A_{k-1}$:
\[
= P(A_1 \cap \dots \cap A_{k-2}) P(A_{k-1} \mid A_1 \cap \dots \cap A_{k-2}) P(A_k \mid A_1 \cap \dots \cap A_{k-1})
\]
I repeat this process recursively until I reach the base case $P(A_1)$, which yields the fully expanded formula:
\[
= P(A_1) P(A_2 \mid A_1) P(A_3 \mid A_1 \cap A_2) \dots P(A_k \mid A_1 \cap A_2 \cap \dots \cap A_{k-1})
\]

\vspace{0.5cm}

\textbf{\large Page 6 of Ch 2-part 2 notes: Use Venn diagram to explain the formula for $P(A \cup B \cup C)$}

To find the total probability of the three overlapping circles $A$, $B$, and $C$:
\begin{itemize}
    \item I add the individual probabilities $P(A) + P(B) + P(C)$.
    \item I subtract the pairwise intersections $P(A \cap B)$, $P(A \cap C)$, and $P(B \cap C)$ because adding the individual circles counted these overlapping regions twice.
    \item I add the central intersection $P(A \cap B \cap C)$ back once. It was initially added three times (in $A$, $B$, and $C$) and subtracted three times (in the pairwise intersections), giving a net count of zero. Adding it back ensures the region is counted exactly once.
\end{itemize}

\vspace{0.5cm}

\textbf{\large Example 7, Page 8 of Ch 2-part 2 notes: Find $P(A) = P(B_1 \cup B_2 \cup B_3)$ without using De Morgan's law}

I know the probability of each independent event is $P(B_1) = P(B_2) = P(B_3) = 0.9$.
I use the additive law for three events:
\[
P(B_1 \cup B_2 \cup B_3) = P(B_1) + P(B_2) + P(B_3) - P(B_1 \cap B_2) - P(B_1 \cap B_3) - P(B_2 \cap B_3) + P(B_1 \cap B_2 \cap B_3)
\]
Because the events are independent, the probability of their intersections is the product of their individual probabilities:
\[
P(B_i \cap B_j) = (0.9)(0.9) = 0.81
\]
\[
P(B_1 \cap B_2 \cap B_3) = (0.9)(0.9)(0.9) = 0.729
\]
I substitute these values into the formula:
\[
P(B_1 \cup B_2 \cup B_3) = 0.9 + 0.9 + 0.9 - 0.81 - 0.81 - 0.81 + 0.729
\]
\[
= 2.7 - 2.43 + 0.729 = 0.27 + 0.729 = 0.999
\]

\end{document}